\par\Needspace{6\baselineskip}\medskip
\begin{center}\hypertarget{native-171-article124}{}\pdfbookmark[2]{124.}{article-124}124.\end{center}
Progredimur iam ad enumerationem methodorum maxime idonearum principiis
praecedentibus innixarum, quarum quidem praecipua momenta in Libro
primo exposita sunt, atque hic tantum instituto nostro accomodari debent.

Methodus simplicissima esse videtur, si pro $x,y$ distantiae corporis coelestis
a terra in duabus observationibus accipiantur, aut potius vel logarithmi
harum distantiarum vel logarithmi distantiarum ad eclipticam sive aequatorem
proiectarum. Hinc per art. 64, V elicientur loca heliocentrica et distantiae a
Sole ad eadem loca pertinentia; hinc porro per art. 110 situs plani orbitae
atque longitudines heliocentricae in ea; hinc atque ex radiis vectoribus temporibusque
respondentibus per problema in artt. 85—105 copiose pertractatum
cuncta reliqua elementa, per quae illas observationes exacte repraesentari
manifestum est, quicunque valores ipsis $x,y$ tributi fuerint. Quodsi iam per
haec elementa locus geocentricus pro tempore observationis tertiae computatur,
huius consensus cum observato vel dissensus decidet, utrum valores suppositi
veri fuerint, an ab iis discrepent; unde quum comparatio duplex derivetur,
differentia altera (in longitudine vel ascensione recta) accipi poterit pro
$X$, alteraque (in latitudine vel declinatione) pro $Y$. Nisi igitur valores harum
differentiarum $X,Y$ sponte prodeunt $=0$, valores veros ipsarum $x,y$ per methodum
in artt. 120 sqq. descriptam eruere licebit. Ceterum per se arbitrarium
est, a quibusnam trium observationum proficiscamur: plerumque tamen praestat,
primam et postremam adoptare, casu speciali de quo statim dicemus excepto.

\SourceLeaf{172}Haecce methodus plerisque post explicandis eo nomine praeferenda est,
quod applicationem maxime generalem patitur. Excipere oportet casum, ubi
duae observationes extremae motum heliocentricum 180 vel 360 vel 540 etc.
graduum complectuntur; tunc enim positio plani orbitae e duobus locis heliocentricis
determinari nequit (art. 110). Perinde methodum applicare haud conveniet,
quoties motus heliocentricus inter duas observationes extremas perparum
differt ab $180^0$ vel $360^0$ etc., quoniam in hoc casu determinatio positionis
orbitae accurata obtineri nequit, sive potius, quoniam variationes levissimae
in valoribus suppositis incognitarum tantas variationes in positione orbitae
et proin etiam in valoribus ipsarum $X,Y$ producerent, ut hae illis non
amplius proportionales censeri possent. Verumtamen remedium hic praesto
est; scilicet in tali casu non proficiscemur a duabus observationibus extremis,
sed a prima et media, vel a media et ultima, adeoque pro $X,Y$ accipiemus
differentias inter computum et observationem in loco tertio vel primo. Quodsi
autem tum locus secundus a primo tum tertius a secundo propemodum 180
gradibus distarent, incommodum illud hoc modo tollere non liceret; sed praestat,
huiusmodi observationes, e quibus per rei naturam determinatio accurata
situs orbitae erui omnino nequit, ad calculum elementorum haud adhibere.

Praeterea haec methodus eo quoque se commendat, quod nullo negotio
aestimari potest, quantas variationes elementa patiantur, si manentibus locis
extremis medius paullulum mutetur: hoc itaque modo iudicium ferri poterit
qualecunque de gradu praecisionis elementis inventis tribuendae.
